% GENERATED FILE. Edit canonical lessons, then run npm run generate:books.
% canonical-source-hash: fnv1a64:c774f73839feca67
% canonical-lessons: TE-C512-jabbu, TE-C512-jabbupadu, TE-C512-nalata, TE-C512-mancam-pattu, TE-C512-teruko

\chapter{Falling Ill}
\label{ch:te-falling-ill}
\hlchaptermodality{512}

\begin{chapteropening}
\textbf{By the end of this chapter:} \emph{I can say that someone fell ill, was laid up or came round.}

Complete the last lesson of chapter 512: I can say that someone fell ill, was laid up or came round.
\end{chapteropening}

\section[\texorpdfstring{\te{జబ్బు} (jabbu)}{జబ్బు (jabbu)}]{\te{జబ్బు} (jabbu) — a sickness; an illness}

\label{lesson:TE-C512-jabbu}



\begin{quote}
Before the new one: say the Telugu for to go over again, then the Telugu for to remind.
\end{quote}

\subsection*{You'll want to know: \te{జబ్బు}}
\textbf{\te{జబ్బు}} — \emph{jabbu} — \textquotedblleft{}a sickness; an illness\textquotedblright{}.

\begin{quote}
\te{ఆయనకు} \te{గుండె} \te{జబ్బు} \te{ఉంది}.
\end{quote}

\begin{quote}
\emph{āyanaku guṇḍe jabbu undi.}
\end{quote}

\begin{quote}
He has heart trouble.
\end{quote}

An everyday word: \te{గుండె} \te{జబ్బు}, heart disease. \te{జబ్బుగా} \te{ఉంది}, (he) is unwell.

\begin{cousinweb}
\textbf{\te{జబ్బు}} (\emph{jabbu}) is native Dravidian: DEDR 2338 lists Telugu \emph{jabbu}, \textquotedblleft{}illness, sickness, weakness\textquotedblright{}, in the article of Tamil \emph{cappai}, \textquotedblleft{}that which is weak\textquotedblright{}.
\end{cousinweb}

\begin{grammarlens}[title={What you've built}]
One more word for describing what you lived and felt.
\end{grammarlens}

\subsection*{Your turn}
\begin{itemize}
\raggedright
  \item \emph{Say it:} \emph{jabbu}
  \item \emph{Say it:} \emph{jabbu}, once more
  \item \emph{Say it:} \emph{āyanaku guṇḍe jabbu undi.} — He has heart trouble.
  \item \emph{Recall:} say the Telugu for to go over again, then the Telugu for to remind, then say \emph{jabbu} again
\end{itemize}

\begin{tcolorbox}[breakable,colback=teal!4,colframe=teal!35!black,title={Before you move on}]
What does \te{జబ్బు} mean? (\textquotedblleft{}A sickness; an illness\textquotedblright{}.) Say it once more.
\end{tcolorbox}
\section[\texorpdfstring{\te{జబ్బుపడు} (jabbupaḍu)}{జబ్బుపడు (jabbupaḍu)}]{\te{జబ్బుపడు} (jabbupaḍu) — to fall ill}

\label{lesson:TE-C512-jabbupadu}



\begin{quote}
Before the new one: say the Telugu for to remind, then the Telugu for a sickness.
\end{quote}

\subsection*{You'll want to know: \te{జబ్బుపడు}}
\textbf{\te{జబ్బుపడు}} — \emph{jabbupaḍu} — \textquotedblleft{}to fall ill\textquotedblright{}.

\begin{quote}
\te{వర్షంలో} \te{తడిసి} \te{తమ్ముడు} \te{జబ్బుపడ్డాడు}.
\end{quote}

\begin{quote}
\emph{varṣaṁlō taḍisi tammuḍu jabbupaḍḍāḍu.}
\end{quote}

\begin{quote}
My younger brother got wet in the rain and fell ill.
\end{quote}

Past (I …): \te{జబ్బుపడ్డాను}, \emph{jabbupaḍḍānu}.

\te{జబ్బు} + \te{పడు}: to fall into illness.

\begin{cousinweb}
\textbf{\te{జబ్బుపడు}} (\emph{jabbupaḍu}) is native Dravidian: \emph{jabbu}, \textquotedblleft{}illness\textquotedblright{} (DEDR 2338), and the native verb \emph{paḍu}, \textquotedblleft{}to feel, to suffer\textquotedblright{} (DEDR 3853).
\end{cousinweb}

\begin{grammarlens}[title={What you've built}]
One more word for describing what you lived and felt.
\end{grammarlens}

\subsection*{Your turn}
\begin{itemize}
\raggedright
  \item \emph{Say it:} \emph{jabbupaḍu}
  \item \emph{Say it:} \emph{jabbupaḍu}, once more
  \item \emph{Say it:} \emph{varṣaṁlō taḍisi tammuḍu jabbupaḍḍāḍu.} — My younger brother got wet in the rain and fell ill.
  \item \emph{Recall:} say the Telugu for to remind, then the Telugu for a sickness, then say \emph{jabbupaḍu} again
\end{itemize}

\begin{tcolorbox}[breakable,colback=teal!4,colframe=teal!35!black,title={Before you move on}]
What does \te{జబ్బుపడు} mean? (\textquotedblleft{}To fall ill\textquotedblright{}.) Say it once more.
\end{tcolorbox}
\section[\texorpdfstring{\te{నలత} (nalata)}{నలత (nalata)}]{\te{నలత} (nalata) — a slight indisposition}

\label{lesson:TE-C512-nalata}



\begin{quote}
Before the new one: say the Telugu for a sickness, then the Telugu for to fall ill.
\end{quote}

\subsection*{You'll want to know: \te{నలత}}
\textbf{\te{నలత}} — \emph{nalata} — \textquotedblleft{}a slight indisposition\textquotedblright{}.

\begin{quote}
\te{నిన్నటి} \te{నుంచి} \te{నాకు} \te{కొంచెం} \te{నలతగా} \te{ఉంది}.
\end{quote}

\begin{quote}
\emph{ninnaṭi nuñci nāku koñceṁ nalatagā undi.}
\end{quote}

\begin{quote}
Since yesterday I have been feeling a bit off.
\end{quote}

Milder than \te{జబ్బు}: \te{నలతగా} \te{ఉంది}, (I) feel a bit off.

\begin{cousinweb}
\textbf{\te{నలత}} (\emph{nalata}) is native Dravidian: DEDR 3611 lists Telugu \emph{nalata}, \textquotedblleft{}disease\textquotedblright{}, beside \emph{naliyu}, \textquotedblleft{}to be crushed, to be bruised\textquotedblright{}, and Tamil \emph{nalaṅku}, \textquotedblleft{}to grow faint\textquotedblright{}.
\end{cousinweb}

\begin{grammarlens}[title={What you've built}]
One more word for describing what you lived and felt.
\end{grammarlens}

\subsection*{Your turn}
\begin{itemize}
\raggedright
  \item \emph{Say it:} \emph{nalata}
  \item \emph{Say it:} \emph{nalata}, once more
  \item \emph{Say it:} \emph{ninnaṭi nuñci nāku koñceṁ nalatagā undi.} — Since yesterday I have been feeling a bit off.
  \item \emph{Recall:} say the Telugu for a sickness, then the Telugu for to fall ill, then say \emph{nalata} again
\end{itemize}

\begin{tcolorbox}[breakable,colback=teal!4,colframe=teal!35!black,title={Before you move on}]
What does \te{నలత} mean? (\textquotedblleft{}A slight indisposition\textquotedblright{}.) Say it once more.
\end{tcolorbox}
\section[\texorpdfstring{\te{మంచం} \te{పట్టు} (mañcaṁ paṭṭu)}{మంచం పట్టు (mañcaṁ paṭṭu)}]{\te{మంచం} \te{పట్టు} (mañcaṁ paṭṭu) — to be bedridden; to take to one's bed}

\label{lesson:TE-C512-mancam-pattu}



\begin{quote}
Before the new one: say the Telugu for to fall ill, then the Telugu for a slight indisposition.
\end{quote}

\subsection*{You'll want to know: \te{మంచం} \te{పట్టు}}
\textbf{\te{మంచం} \te{పట్టు}} — \emph{mañcaṁ paṭṭu} — \textquotedblleft{}to be bedridden; to take to one's bed\textquotedblright{}.

\begin{quote}
\te{జ్వరంతో} \te{తాతయ్య} \te{వారం} \te{రోజులు} \te{మంచం} \te{పట్టారు}.
\end{quote}

\begin{quote}
\emph{jvarantō tātayya vāraṁ rōjulu mañcaṁ paṭṭāru.}
\end{quote}

\begin{quote}
With fever, grandpa was in bed for a week.
\end{quote}

Past (I …): \te{మంచం} \te{పట్టాను}, \emph{mañcaṁ paṭṭānu}.

Literally to take hold of the bed: \te{మంచం}, bed, + \te{పట్టు}, to hold.

\begin{cousinweb}
\textbf{\te{మంచం} \te{పట్టు}} (\emph{mañcaṁ paṭṭu}) is a Telugu phrase: \te{మంచం}, \textquotedblleft{}a bed\textquotedblright{}, from Sanskrit \emph{mañca}, \textquotedblleft{}a platform, a bed\textquotedblright{} (CDIAL 9715), and the native verb \emph{paṭṭu}, \textquotedblleft{}to hold, to seize\textquotedblright{} (DEDR 4034).
\end{cousinweb}

\begin{grammarlens}[title={What you've built}]
One more word for describing what you lived and felt.
\end{grammarlens}

\subsection*{Your turn}
\begin{itemize}
\raggedright
  \item \emph{Say it:} \emph{mañcaṁ paṭṭu}
  \item \emph{Say it:} \emph{mañcaṁ paṭṭu}, once more
  \item \emph{Say it:} \emph{jvarantō tātayya vāraṁ rōjulu mañcaṁ paṭṭāru.} — With fever, grandpa was in bed for a week.
  \item \emph{Recall:} say the Telugu for to fall ill, then the Telugu for a slight indisposition, then say \emph{mañcaṁ paṭṭu} again
\end{itemize}

\begin{tcolorbox}[breakable,colback=teal!4,colframe=teal!35!black,title={Before you move on}]
What does \te{మంచం} \te{పట్టు} mean? (\textquotedblleft{}To be bedridden; to take to one's bed\textquotedblright{}.) Say it once more.
\end{tcolorbox}
\section[\texorpdfstring{\te{తేరుకో} (tērukō)}{తేరుకో (tērukō)}]{\te{తేరుకో} (tērukō) — to come round; to recover from a shock}

\label{lesson:TE-C512-teruko}



\begin{quote}
Before the new one: say the Telugu for a slight indisposition, then the Telugu for to be bedridden.
\end{quote}

\subsection*{You'll want to know: \te{తేరుకో}}
\textbf{\te{తేరుకో}} — \emph{tērukō} — \textquotedblleft{}to come round; to recover from a shock\textquotedblright{}.

\begin{quote}
\te{ఆ} \te{వార్త} \te{నుంచి} \te{అమ్మ} \te{తేరుకోవడానికి} \te{వారం} \te{పట్టింది}.
\end{quote}

\begin{quote}
\emph{ā vārta nuñci amma tērukōvaḍāniki vāraṁ paṭṭindi.}
\end{quote}

\begin{quote}
It took mum a week to get over that news.
\end{quote}

Past (I …): \te{తేరుకున్నాను}, \emph{tērukunnānu}.

-\te{నుంచి} \te{తేరుకో}, to get over a shock, a fright or a faint.

\begin{cousinweb}
\textbf{\te{తేరుకో}} (\emph{tērukō}) is native Dravidian: DEDR 3471 lists Telugu \emph{tēṟu} with Tamil \emph{tēṟu}, \textquotedblleft{}to recover from swooning, to be comforted, to cheer up\textquotedblright{}; -\te{కో} is the reflexive ending.
\end{cousinweb}

\begin{grammarlens}[title={What you've built}]
One more word for describing what you lived and felt.
\end{grammarlens}

\subsection*{Your turn}
\begin{itemize}
\raggedright
  \item \emph{Say it:} \emph{tērukō}
  \item \emph{Say it:} \emph{tērukō}, once more
  \item \emph{Say it:} \emph{ā vārta nuñci amma tērukōvaḍāniki vāraṁ paṭṭindi.} — It took mum a week to get over that news.
  \item \emph{Recall:} say the Telugu for a slight indisposition, then the Telugu for to be bedridden, then say \emph{tērukō} again
\end{itemize}

\begin{tcolorbox}[breakable,colback=teal!4,colframe=teal!35!black,title={Before you move on}]
What does \te{తేరుకో} mean? (\textquotedblleft{}To come round; to recover from a shock\textquotedblright{}.) Say it once more.
\end{tcolorbox}
